% Template autoral SquadOS — MIT. Pacotes LaTeX vêm da distribuição TeX.
\documentclass[11pt,a4paper]{article}
\usepackage[T1]{fontenc}
\usepackage[utf8]{inputenc}
\usepackage{graphicx}
\usepackage{booktabs}
\usepackage{hyperref}
\title{Título da pesquisa}
\author{Autoria a confirmar}
\date{}
\begin{document}
\maketitle
\begin{abstract}
% Pergunta, método, resultado observado e limites. Não inventar resultados.
Resumo a elaborar com base no trabalho realizado.
\end{abstract}
\section{Introdução}
% Contexto, lacuna demonstrada e pergunta de pesquisa.
\section{Método}
% Dados, seleção, procedimento reproduzível e critérios de análise.
\section{Resultados}
% Achados observados, unidades, incerteza e fontes dos dados.
\section{Discussão e limitações}
% Interpretação, alternativas e limites de generalização.
\section{Conclusão}
% Responder à pergunta de modo proporcional à evidência.
\begin{thebibliography}{9}
% Adicionar apenas referências reais conferidas nas fontes consultadas.
\end{thebibliography}
\end{document}
